\section{CUDA 的历史性决策：把并行能力从“图形技巧”变成“通用基础设施”}

\subsection{从“把科学问题伪装成图形问题”到标准化接口}

Jensen 回忆了早期 GPGPU（General-purpose computing on GPU）阶段的尴尬：研究者必须“欺骗”图形流水线，让 GPU 误以为自己在做图形任务，才能借用其并行吞吐。这个路径能跑通实验，但无法形成工程生态。

CUDA 的决定性价值在于：把 GPU 编程从图形 API 语义中解耦，转为通用语言和通用开发工具链。访谈里他强调“你可以用 C 去告诉 GPU 要做什么”，本质上就是在做“开发者生产函数”的提升。

\begin{figure}[H]
\centering
\includegraphics[width=0.88\textwidth]{frame-cuda.jpg}
\caption{CUDA 作为软件层，把 GPU Program Execution 对外开放 \protect\footnotemark}
\end{figure}
\footnotetext{视频画面时间区间：00:08:29--00:08:40。}

\begin{importantbox}{If you do not build it they cannot come}
Jensen 给出平台投资的核心逻辑：\textit{``If you build it, they might not come; but if you don't build it, they can't come.''}

这是典型的先验平台投资：先承受不确定性，再换取生态上限。CUDA 在诞生时并没有“确定的大客户合同”，但它提前铺设了后来 AI 爆发所需的软件底座。
\end{importantbox}

\subsection{为什么 NVIDIA 敢押这么大：量产市场提供“经济可行性”}

访谈里一个常被忽略但极关键的点是：Jensen 并非“盲目理想主义”。他反复提到“视频游戏市场的规模”，因为这意味着 GPU 会成为全球高产量并行处理器。高产量带来的不是营销声量，而是两个硬约束：
\begin{itemize}
    \item 成本曲线可下探，研究机构和创业团队可负担。
    \item 迭代节奏可维持，软件生态不至于断代。
\end{itemize}

\begin{warningbox}{误区：CUDA 的价值只在“性能提升”}
如果只把 CUDA 理解为“让程序跑快一点”，会低估它的历史意义。CUDA 的真正价值是把并行计算从“高手手工活”变成“可教学、可复用、可规模化的工业能力”。这类能力转移通常比单次性能提升更有长期复利。
\end{warningbox}

\subsection{本章小结}
CUDA 不是一个工具发布，而是一次平台级组织承诺。它把 GPU 从“图形芯片”重定义为“通用并行计算底座”，为后续 AI 产业化提供了先决条件。

